%================================================================
\section{The analysts: a ladder of assumptions}\label{sec:analysts}

Each analyst sees one panel at a diagnostic point $c_0$: the outcomes $Y_T$, the
assignment, the covariates and the graph, from one draw of the simulation noise.
None of them sees NROY. Three analysts report an interval for $\theta_2$; one,
Foil~2, reports an interval for a different contrast, the spillover on the pilot.
Every interval has nominal level $\ConfidencePct\%$ and uses the two-sided normal
quantile. Table~\ref{tab:ladder} lists them in order of decreasing assumption
strength. A0, A1 and A2 are three intervals aimed at $\theta_2$ under different
assumptions. They do not estimate the same functional, and only A0's needs
SUTVA to equal $\theta_2$.

\begin{table}[t]
\centering
\caption{The analysts. A0, A1 and A2 aim at $\theta_2$; Foil~2 aims at the
spillover $\dspill$ and is compared only with its own image.}
\label{tab:ladder}
\begin{zebratabularx}{\linewidth}{lllX}
\toprule
\colhead{Analyst} & \colhead{Panel} & \colhead{Target} & \colhead{Maintained assumptions} \\
\midrule
A0 (Foil 1) & policy arm & $\theta_2$ & no interference; unconfoundedness given $X$; linear constant-effect model \\
A1 & pilot & $\theta_2$ & exposure map; transport from the pilot with the policy's exposure shares \\
A2 & pilot & $\theta_2$ & exposure map; transport within the hubs and within everyone else \\
Foil 2 & pilot & $\dspill$ & exposure map; its target is a spillover contrast \\
\bottomrule
\end{zebratabularx}
\tablenote{A0 is compared with $\Theta_{\hm}$, as are A1 and A2; Foil~2 is compared
with $\Theta_{\expset}$. A0 and Foil~2 are the two foils of the main
pre-registration; A1 and A2 were added by its amendment A2.}
\end{table}

\subsection{A0: the analyst who assumes no interference (Foil 1)}\label{sec:foil1}

A0 has one panel from the policy itself, the arm that seeds the hubs. She fits OLS
of $Y_T$ on $Z_i=(1,d_i,X_i^\top)$ and reads the coefficient on $d_i$,
$\hat\beta$, with the heteroskedasticity-robust HC1 standard error
\citep{mackinnon1985},
\begin{equation}\label{eq:hc1}
  \widehat{\Var}_{\mathrm{HC1}}(\hat\beta)
  =\Bigl[\tfrac{n}{n-p}\,(Z^\top Z)^{-1}
   \Bigl(\textstyle\sum_{i}\hat u_i^2\,Z_iZ_i^\top\Bigr)(Z^\top Z)^{-1}\Bigr]_{dd},
\end{equation}
where $p$ is the number of columns of $Z$. Under no interference, the policy effect
for the whole village is the seeded share times the effect of own seeding, so she
reports $(K/N)\hat\beta$ with the interval scaled the same way, and we call that
interval $\Jone$. The scaling by a known positive constant keeps her own coverage
statement and does not repair her assumptions.

Her interval answers the question it was built for under three conditions: no
interference, so that a person's outcome depends only on her own seeding;
unconfoundedness given $X$; and a linear, constant-effect outcome model. The
village breaks the first by construction. The second is not supported either: on
the policy arm, seeding is a deterministic function of degree, which the
regression does not include. A0 is still the natural default, and it is a good
estimator of what it estimates. Under interference, $(K/N)\hat\beta$ measures
something close to the effect of being one of the hubs on the hubs, and it has no
term for the effect of the campaign on everyone else. Appendix~\ref{app:gap}
decomposes the gap exactly and explains why a flexible regression on the same
panel would not help.

\subsection{Foil 2: the spillover on the pilot}\label{sec:foil2}

Foil~2 uses the pilot panel and the exposure strata of
Section~\ref{sec:assumptions}. Because the pilot's lottery is known, the
probability $p_{i,S}$ that person $i$ lands in stratum $S$ can be computed before
any outcome is seen. For an unseeded person of degree $k$, the number of seeded
neighbors under the pilot is hypergeometric, so the first-order probabilities are
exact. Foil~2 replaces each stratum mean with a H\'ajek mean,
\begin{equation}\label{eq:hajek}
  \widehat\mu_S
  =\frac{\sum_i \mathbf{1}\{i\in S\}\,Y_{iT}/p_{i,S}}{\sum_i \mathbf{1}\{i\in S\}/p_{i,S}},
\end{equation}
and estimates the spillover contrast
$\widehat\mu_{\Sc}-\widehat\mu_{\Sa}$: the outcome of unseeded people with many
seeded neighbors against unseeded people with none. Its target is
\begin{equation}\label{eq:spill}
  \dspill(c)=\overline Y_{\Sc}-\overline Y_{\Sa},
\end{equation}
the difference of the realized stratum means of the true outcomes, averaged over
$\RepsEstimand$ pilot replications. The variance is the linearized estimator of
\citet{aronow2017}. Its joint exposure probabilities come from $\ExposureDraws$
simulated pilot assignments, and Young's bound, as used by \citet{aronow2017},
covers pairs of people who never appear together in a pair of strata. If any
stratum in a panel holds fewer than $\StratumMinN$ people, the code merges $\Sb$
into $\Sc$ for that panel. We call the interval $\Jtwo$.

Foil~2 is the analyst who takes interference seriously but answers a different
question. Its target is $\dspill$, not $\theta_2$, and its interval is compared
only with $\Theta_{\expset}$, the image of $\dspill$. The H\'ajek mean targets an
inverse-probability-weighted exposure contrast, while $\dspill$ averages the
realized contrasts without weights; on this run the two are close enough that no
containment flag depends on the difference (Appendix~\ref{app:foil2}). Foil~2 is
the only analyst with a pre-registered calibration check of its interval, also in
Appendix~\ref{app:foil2}.

\subsection{A1: exposure transport}\label{sec:a1}

A1 uses the same pilot panel, the same strata, the same cut and the same merge
rule as Foil~2, and the same H\'ajek means with exact first-order probabilities.
It estimates Equation~\eqref{eq:transport},
\begin{equation}\label{eq:a1}
  \widehat\theta_2^{\,\mathrm{A1}}
  =\sum_{S\neq\Sa} w_S\bigl(\widehat\mu_S-\widehat\mu_{\Sa}\bigr),
  \qquad w_S=\frac{n_S(\dtwo)}{N},
\end{equation}
where $n_S(\dtwo)$ is the number of people in stratum $S$ under the policy vector.
When a panel is merged, the shares use the merged strata. The estimate is a fixed
linear combination of H\'ajek means with coefficients $a_S$, so its variance is
the same linearized covariance summed over all pairs of strata,
$\sum_S\sum_{S'}a_S a_{S'}\,\widehat{\Cov}(\widehat\mu_S,\widehat\mu_{S'})$, with
joint probabilities from the same simulated assignments.

A1 maintains the exposure map and transport. It allows interference, but only
the part that runs through the share of seeded neighbors, and it applies the
pilot's exposure-specific outcomes to the people who would hold those exposures
under the policy.

\subsection{A2: degree-binned transport}\label{sec:a2}

A2 splits the village into two bins fixed by the graph: $H$, the $K$ people whom
$\dtwo$ seeds, and $L$, everyone else. Under $\dtwo$ every person in $H$ is in
$\Sone$, so
\begin{equation}\label{eq:a2}
  \widehat\theta_2^{\,\mathrm{A2}}
  =\frac{K}{N}\bigl(\widehat\mu_{\Sone,H}-\widehat\mu_{\Sa,H}\bigr)
   +\sum_{S\neq\Sa}\frac{n_{S,L}(\dtwo)}{N}\bigl(\widehat\mu_{S,L}-\widehat\mu_{\Sa,L}\bigr),
\end{equation}
where $\widehat\mu_{S,B}$ is the H\'ajek mean of stratum $S$ within bin $B$. The
variance uses the same machinery, with each cell's indicators restricted to its
bin. A2 never merges across bins and never borrows a control mean from the other
bin. If a cell that Equation~\eqref{eq:a2} needs is empty in a panel, A2 is
undefined for that panel; the code counts that case and never imputes it.

A2 is the weakest rung of the ladder. The hub term compares the hubs with
themselves, so it needs the exposure map but no transport across people. The
remaining terms transport only within $L$. The price is precision: A2 restricts
each cell to one bin, so the cells it needs are thinner than those of A1, which
pools them.

\subsection{What the ladder predicts, written in advance}\label{sec:expectations}

Amendment A2 of the pre-registration wrote five expectations before any A2
number existed. At every diagnostic point:
(a)~the absolute bias falls from A0 to A1 to A2;
(b)~the sampling SD and the interval width rise from A0 to A1 to A2;
(c)~A0's coverage is far below nominal, and A2's is the closest to nominal of the
three;
(d)~A1 and A2 stay biased downward, because two-hop spillover raises the people
in $\Sa$ above the true status-quo level;
(e)~every interval is narrower than $\Theta_{\hm}$.
Any of them could fail, and a failure is reported as found.
Section~\ref{sec:a2results} checks each one.
